% GENERATED FILE. Edit canonical lessons, then run npm run generate:books.
% canonical-source-hash: fnv1a64:4c6a18fa1442e1b7
% canonical-lessons: TA-C140-tiratcai, TA-C140-matulai, TA-C140-urukay, TA-C140-aruvi, TA-C140-kunru

\chapter{Grapes, Pomegranate, Pickle, Waterfall, Hill}
\label{ch:ta-grapes-pomegranate-pickle-waterfall-hill}
\hlchaptermodality{140}

\begin{chapteropening}
\textbf{By the end of this chapter:} \emph{I can say grapes, a pomegranate, a pickle, a waterfall and a hill.}

Complete the last lesson of chapter 140: I can say grapes, a pomegranate, a pickle, a waterfall and a hill.
\end{chapteropening}

\section[\texorpdfstring{\ta{திராட்சை} (tirāṭcai)}{திராட்சை (tirāṭcai)}]{\ta{திராட்சை} (tirāṭcai) — grapes}

\label{lesson:TA-C140-tiratcai}



\begin{quote}
Before the new one: say the Tamil for a guava, then the Tamil for a lemon.
\end{quote}

\subsection*{You'll want to know: \ta{திராட்சை}}
\textbf{\ta{திராட்சை}} — \emph{tirāṭcai} — \textquotedblleft{}grapes\textquotedblright{}.

\textbf{In the family, and next door}

\noindent
\begin{tabularx}{\linewidth}{@{}>{\raggedright\arraybackslash}X>{\raggedright\arraybackslash}X@{}}
\toprule
\textbf{} & \textbf{word} \\
\midrule
Kannada & \ka{ದ್ರಾಕ್ಷಿ} (\emph{drākṣi}) \\
Telugu & \te{ద్రాక్ష} (\emph{drākṣa}) \\
Malayalam & \ml{മുന്തിരി} (\emph{muntiri}) \\
Hindi (neighbour) & \dv{अंगूर} (\emph{aṅgūr}) \\
English & grapes \\
\bottomrule
\end{tabularx}

\begin{grammarlens}[title={What you've built}]
One more word for this chapter.
\end{grammarlens}

\subsection*{Your turn}
\begin{itemize}
\raggedright
  \item \emph{Say it:} \emph{tirāṭcai}
  \item \emph{Say it:} \emph{tirāṭcai}, once more
  \item \emph{Say it:} say \emph{tirāṭcai}
  \item \emph{Recall:} say the Tamil for a guava, then the Tamil for a lemon, then say \emph{tirāṭcai} again
\end{itemize}

\begin{tcolorbox}[breakable,colback=teal!4,colframe=teal!35!black,title={Before you move on}]
What does \ta{திராட்சை} mean? (\textquotedblleft{}Grapes\textquotedblright{}.) Say it once more.
\end{tcolorbox}
\section[\texorpdfstring{\ta{மாதுளை} (mātuḷai)}{மாதுளை (mātuḷai)}]{\ta{மாதுளை} (mātuḷai) — a pomegranate}

\label{lesson:TA-C140-matulai}



\begin{quote}
Before the new one: say the Tamil for a lemon, then the Tamil for grapes.
\end{quote}

\subsection*{You'll want to know: \ta{மாதுளை}}
\textbf{\ta{மாதுளை}} — \emph{mātuḷai} — \textquotedblleft{}a pomegranate\textquotedblright{}.

\begin{grammarlens}[title={What you've built}]
One more word for this chapter.
\end{grammarlens}

\subsection*{Your turn}
\begin{itemize}
\raggedright
  \item \emph{Say it:} \emph{mātuḷai}
  \item \emph{Say it:} \emph{mātuḷai}, once more
  \item \emph{Say it:} say \emph{mātuḷai}
  \item \emph{Recall:} say the Tamil for a lemon, then the Tamil for grapes, then say \emph{mātuḷai} again
\end{itemize}

\begin{tcolorbox}[breakable,colback=teal!4,colframe=teal!35!black,title={Before you move on}]
What does \ta{மாதுளை} mean? (\textquotedblleft{}A pomegranate\textquotedblright{}.) Say it once more.
\end{tcolorbox}
\section[\texorpdfstring{\ta{ஊறுகாய்} (ūṟukāy)}{ஊறுகாய் (ūṟukāy)}]{\ta{ஊறுகாய்} (ūṟukāy) — a pickle}

\label{lesson:TA-C140-urukay}



\begin{quote}
Before the new one: say the Tamil for grapes, then the Tamil for a pomegranate.
\end{quote}

\subsection*{You'll want to know: \ta{ஊறுகாய்}}
\textbf{\ta{ஊறுகாய்}} — \emph{ūṟukāy} — \textquotedblleft{}a pickle\textquotedblright{}.

\textbf{In the family, and next door}

\noindent
\begin{tabularx}{\linewidth}{@{}>{\raggedright\arraybackslash}X>{\raggedright\arraybackslash}X@{}}
\toprule
\textbf{} & \textbf{word} \\
\midrule
Kannada & \ka{ಉಪ್ಪಿನಕಾಯಿ} (\emph{uppinakāyi}) \\
English & a pickle \\
\bottomrule
\end{tabularx}

\begin{grammarlens}[title={What you've built}]
One more word for this chapter.
\end{grammarlens}

\subsection*{Your turn}
\begin{itemize}
\raggedright
  \item \emph{Say it:} \emph{ūṟukāy}
  \item \emph{Say it:} \emph{ūṟukāy}, once more
  \item \emph{Say it:} say \emph{ūṟukāy}
  \item \emph{Recall:} say the Tamil for grapes, then the Tamil for a pomegranate, then say \emph{ūṟukāy} again
\end{itemize}

\begin{tcolorbox}[breakable,colback=teal!4,colframe=teal!35!black,title={Before you move on}]
What does \ta{ஊறுகாய்} mean? (\textquotedblleft{}A pickle\textquotedblright{}.) Say it once more.
\end{tcolorbox}
\section[\texorpdfstring{\ta{அருவி} (aruvi)}{அருவி (aruvi)}]{\ta{அருவி} (aruvi) — a waterfall}

\label{lesson:TA-C140-aruvi}



\begin{quote}
Before the new one: say the Tamil for a pomegranate, then the Tamil for a pickle.

Say \textbf{\ta{வானிலை}}, then \textbf{\ta{மழை} \ta{பெய்கிறது}}. Name the rainy season, \textbf{\ta{மழைக்} \ta{காலம்}}, the one that shapes the year here.
\end{quote}

\subsection*{You'll want to know: \ta{அருவி}}
\textbf{\ta{அருவி}} — \emph{aruvi} — \textquotedblleft{}a waterfall\textquotedblright{}.

\textbf{In the family, and next door}

\noindent
\begin{tabularx}{\linewidth}{@{}>{\raggedright\arraybackslash}X>{\raggedright\arraybackslash}X@{}}
\toprule
\textbf{} & \textbf{word} \\
\midrule
Kannada & \ka{ಜಲಪಾತ} (\emph{jalapāta}) \\
Telugu & \te{జలపాతం} (\emph{jalapātaṁ}) \\
Malayalam & \ml{വെള്ളച്ചാട്ടം} (\emph{veḷḷaccāṭṭaṁ}) \\
English & a waterfall \\
\bottomrule
\end{tabularx}

\begin{grammarlens}[title={What you've built}]
One more word for this chapter.
\end{grammarlens}

\subsection*{Your turn}
\begin{itemize}
\raggedright
  \item \emph{Say it:} \emph{aruvi}
  \item \emph{Say it:} \emph{aruvi}, once more
  \item \emph{Say it:} say \emph{aruvi}
  \item \emph{Recall:} say the Tamil for a pomegranate, then the Tamil for a pickle, then say \emph{aruvi} again
\end{itemize}

\begin{tcolorbox}[breakable,colback=teal!4,colframe=teal!35!black,title={Before you move on}]
What does \ta{அருவி} mean? (\textquotedblleft{}A waterfall\textquotedblright{}.) Say it once more.
\end{tcolorbox}
\section[\texorpdfstring{\ta{குன்று} (kuṉṟu)}{குன்று (kuṉṟu)}]{\ta{குன்று} (kuṉṟu) — a hill}

\label{lesson:TA-C140-kunru}



\begin{quote}
Before the new one: say the Tamil for a pickle, then the Tamil for a waterfall.
\end{quote}

\subsection*{You'll want to know: \ta{குன்று}}
\textbf{\ta{குன்று}} — \emph{kuṉṟu} — \textquotedblleft{}a hill\textquotedblright{}.

\textbf{In the family, and next door}

\noindent
\begin{tabularx}{\linewidth}{@{}>{\raggedright\arraybackslash}X>{\raggedright\arraybackslash}X>{\raggedright\arraybackslash}X@{}}
\toprule
\textbf{} & \textbf{word} & \textbf{same root} \\
\midrule
Kannada & \ka{ಬೆಟ್ಟ} (\emph{beṭṭa}) &  \\
Telugu & \te{కొండ} (\emph{koṇḍa}) & yes \\
English & a hill &  \\
\bottomrule
\end{tabularx}

\begin{grammarlens}[title={What you've built}]
One more word for this chapter.
\end{grammarlens}

\subsection*{Your turn}
\begin{itemize}
\raggedright
  \item \emph{Say it:} \emph{kuṉṟu}
  \item \emph{Say it:} \emph{kuṉṟu}, once more
  \item \emph{Say it:} say \emph{kuṉṟu}
  \item \emph{Recall:} say the Tamil for a pickle, then the Tamil for a waterfall, then say \emph{kuṉṟu} again
\end{itemize}

\begin{tcolorbox}[breakable,colback=teal!4,colframe=teal!35!black,title={Before you move on}]
What does \ta{குன்று} mean? (\textquotedblleft{}A hill\textquotedblright{}.) Say it once more.
\end{tcolorbox}
